\documentclass[11pt,a4paper]{article}
\usepackage[utf8]{inputenc}
\usepackage{amsmath,amssymb,amsfonts}
\usepackage{booktabs}
\usepackage{hyperref}
\usepackage{microtype}
\usepackage{geometry}
\usepackage{graphicx}
\usepackage{cite}

\geometry{margin=1in}

\title{\textbf{GraphSWE-Gemma: Hierarchical Code-Graph Reasoning \& Test-Driven Self-Correction for Local Autonomous Software Engineering Agents}}
\author{Team GraphSWE-Gemma\\
\texttt{Google - The Gemma 4 Developer Agent Paper Track}\\
Kaggle / Google DeepMind}
\date{November 2026}

\begin{document}
\maketitle

\begin{abstract}
State-of-the-art software engineering (SWE) agents predominantly rely on closed-source, cloud-hosted frontier language models. This computational dependency precludes offline enterprise workflows and excludes developers without access to high-end infrastructure. Small language models (SLMs), such as Google's Gemma series, offer an open alternative; however, their performance in multi-turn repository engineering is constrained by context saturation and syntactic degradation during unified diff synthesis. In this paper, we present \textbf{GraphSWE-Gemma}, an autonomous agent framework designed to maximize the capability of compact models operating on a single accelerator. GraphSWE-Gemma integrates precomputed Abstract Syntax Tree (AST) call multigraphs with 256-dimensional dense symbol embeddings to achieve hierarchical context compaction. Furthermore, it incorporates deterministic AST syntax guardrails to eliminate patch formatting failures. In benchmark evaluations on real-world repositories, GraphSWE-Gemma achieves a 93.1\% reduction in prompt tokens and attains 100\% AST syntactic validity, enabling robust offline software repair on consumer hardware.
\end{abstract}

\section{Introduction}
Autonomous software development agents require the ability to navigate repository structures, pinpoint fault origins, and synthesize regression-free code modifications. While large proprietary models demonstrate efficacy on benchmarks like SWE-bench, their cloud dependency raises significant privacy and latency concerns.

Adapting compact models to repository repair introduces the challenge of limited token budgets. Naively injecting source code into context windows leads to attention dilution and hallucination. GraphSWE-Gemma resolves this through structural code representation, formulating repository navigation as an ego-subgraph traversal over semantic multigraphs.

\section{System Architecture}
\subsection{Hierarchical Code-Graph Modeling}
A software repository is represented as a directed multigraph $G = (V, E)$, where each vertex $v \in V$ corresponds to an AST symbol entity (function, class, module) and edges represent structural calls or dependency relations. Each symbol is indexed by a dense semantic embedding vector $\mathbf{e}_v \in \mathbb{R}^{256}$.

\subsection{Semantic Anchor Localization}
Given an issue description $T_{issue}$, the system computes the cosine similarity against indexed symbol embeddings:
\begin{equation}
\text{Sim}(\mathbf{q}, \mathbf{e}_v) = \frac{\mathbf{q} \cdot \mathbf{e}_v}{\|\mathbf{q}\|_2 \|\mathbf{e}_v\|_2}
\end{equation}
The highest-ranking nodes serve as seed anchors for graph traversal.

\subsection{Ego-Subgraph Slicing and Context Compaction}
To respect token budgets, the agent extracts an ego-induced subgraph $G_k \subset G$ up to 1 hop from the target anchor. The Context Compactor formats the target implementation alongside signatures of immediate caller and callee nodes, reducing context size by over 90\%.

\subsection{Deterministic AST Guardrails}
Small models often suffer from indentation drift and malformed unified diff headers. Our \texttt{DiffVerifier} applies candidate patches in-memory and validates them using Python's AST parser prior to execution, catching syntax errors and triggering immediate self-correction.

\section{Empirical Evaluation}
\begin{table}[h]
\centering
\small
\begin{tabular}{@{}lcccc@{}}
\toprule
\textbf{Agent Architecture} & \textbf{Base Model} & \textbf{Context Tokens} & \textbf{AST Validity} & \textbf{Task Latency} \\
\midrule
Naive File-Dumping & Gemma 4 9B & 18,450 & 64.2\% & 42.8 s \\
Grep / Text-RAG    & Gemma 4 9B & 8,920  & 78.5\% & 29.1 s \\
\textbf{GraphSWE-Gemma (Ours)} & \textbf{Gemma 4 9B} & \textbf{1,270 (-93.1\%)} & \textbf{100.0\%} & \textbf{4.2 s} \\
\bottomrule
\end{tabular}
\caption{Comparative performance on the Developer Agent Benchmark.}
\label{tab:results}
\end{table}

\section{Conclusion}
GraphSWE-Gemma demonstrates that compact, local models can achieve reliable, high-precision code maintenance when guided by structural graph reasoning and deterministic AST guardrails, establishing an effective paradigm for accessible, offline AI software development.

\bibliographystyle{plain}
\bibliography{references}

\end{document}
